\section{Conclusion}

We studied how institutional processing delay affects the stability of a multi-agent system in which agents adapt their behavior in response to delayed punishment signals. We derived a closed-form critical delay $\Delta_c$ for the delayed replicator equation with a sigmoid response function and proved that the resulting Hopf bifurcation is supercritical for the entire admissible parameter class. We then tested three agent architectures (fixed-policy, reactive, and Q-learning) in a networked simulation with 240 agents across 50 seeds per condition.

The central result answers question~3: learning does not amplify delay-induced instability---it partially buffers it. Non-reactive agents are immune to delay (0\% runaway), reactive agents collapse catastrophically (96\%), and Q-learning agents achieve partial resilience (66\%). The delay sweep (question~1) confirms a monotonic dose-response reaching $+62$ percentage points of excess runaway, the sharpness sweep (question~2) confirms that sharper institutional responses lower the stability threshold, and the RL regulator experiment (question~4) suggests that adaptive governance converts runaway into bounded oscillations rather than restoring stability. The destabilizing ingredient, then, is memoryless reactivity to delayed signals, not learning: agents that immediately exploit low-alarm windows trigger oscillatory feedback loops, while agents with cumulative memory resist this trap. The practical implication is twofold. Institutions should prefer gradual responses over sharp thresholds, because sharpness amplifies the instability that reactivity exploits; and adaptive memory is not the disease but the closest thing to a cure this system has.

Three extensions are natural: scaling to larger populations and alternative network families to test how far the architecture hierarchy holds; continuous action spaces, which may produce different stability characteristics; and combining delay with noisy observation, which the companion paper addresses for selective governance on modular networks.
